\section{Reflected fibers and a sparse incidence theorem}
\label{sec:tpc29-reflected-incidence}

\subsection{Hyperbola reflection}

For a positive target $N_m(j)=mj+h$, divisor reflection in
\eqref{eq:tpc29-ultra-increment} gives the exact identity
\begin{equation}
 C_m(j)=
 \sum_{\substack{r\mid N_m(j)\\r<N_m(j)/T}}
 a\!\left(\frac{N_m(j)}r\right).
 \label{eq:tpc29-reflected-ultra}
\end{equation}
The reflected variable $r$ is not required to be squarefree.  Only the
quotient coefficient $a(N_m/r)$ vanishes off squarefree quotients.

The both-ultra channel contains two variables of the form
\eqref{eq:tpc29-reflected-ultra}.  In a mixed channel, one variable is
an original prefix divisor $u\le T$ and the other is a reflected
complement variable.  Both situations reduce to the same pair of
divisibility conditions
\begin{equation}
 r\mid mj+h,\qquad s\mid nj+h.
 \label{eq:tpc29-two-divisibilities}
\end{equation}
Support restrictions and coefficient types will be restored only
after the common incidence count is proved.

\begin{lemma}[Compatibility and the reflected determinant]
\label{lem:tpc29-reflected-crt}
Let $m\ne n$ be primitive physical rows, and let $r,s$ satisfy
\eqref{eq:tpc29-two-divisibilities} for some primitive orbit point.
Put
\begin{equation}
 d=(r,s),\qquad q=[r,s]=\frac{rs}{d}.
 \label{eq:tpc29-dq}
\end{equation}
Then the two congruences in \eqref{eq:tpc29-two-divisibilities} are
compatible if and only if
\begin{equation}
 d\mid h(m-n).
 \label{eq:tpc29-compatibility}
\end{equation}
When compatible, they select a unique class $j_0\pmod q$.  Writing
$j=j_0+qk$, the two quotient forms
\begin{align}
 U(k)&=\frac{m(j_0+qk)+h}{r}=A_1k+B_1,
 \label{eq:tpc29-U-form}\\
 V(k)&=\frac{n(j_0+qk)+h}{s}=A_2k+B_2
 \label{eq:tpc29-V-form}
\end{align}
have integer coefficients and satisfy
\begin{equation}
 \boxed{A_1B_2-A_2B_1=\frac{h(m-n)}d\ne0.}
 \label{eq:tpc29-affine-determinant}
\end{equation}
Moreover $(d,h)=1$, so compatibility implies $d\mid m-n$.
\end{lemma}

\begin{proof}
Target primitivity gives $(r,m)=(s,n)=1$.  Hence the two congruences
select
\[
 j\equiv-h\overline m\pmod r,
 \qquad
 j\equiv-h\overline n\pmod s.
\]
Their reductions modulo $d$ agree exactly when
$d\mid h(m-n)$.  The generalized Chinese remainder theorem then gives
one class modulo $q$.

The affine coefficients are
\[
 A_1=\frac{mq}{r},\quad B_1=\frac{mj_0+h}{r},
 \qquad
 A_2=\frac{nq}{s},\quad B_2=\frac{nj_0+h}{s}.
\]
Their determinant equals
\[
 \frac{qh(m-n)}{rs}=\frac{h(m-n)}d.
\]
It is nonzero by the off-diagonal mask.  Finally, $d$ divides both
positive targets.  Each target is coprime to $H=\rad|h|$, so
$(d,h)=1$; cancel $h$ in \eqref{eq:tpc29-compatibility}.
\end{proof}

Equivalently, the quotient variables lie on the exact surface
\begin{equation}
 n\,rU-m\,sV=h(n-m).
 \label{eq:tpc29-reflected-surface}
\end{equation}
No coprimality between $d$ and either quotient is asserted.  Such an
assumption would be false for nonsquarefree targets and is not used
below.

\subsection{A dyadic count with multiplicity}

Let $\cI\subset\Z$ be any set contained in an interval of at most
$C_{\mathscr D}J$ consecutive integers, and suppose throughout this
section that
\begin{equation}
 1\le mj+h,\ nj+h\le X^{C_{\mathscr D}}
 \qquad(j\in\cI).
 \label{eq:tpc29-polynomial-target-window}
\end{equation}
The physical target window \eqref{eq:tpc29-target-support} is the case
used later.  For $A,B,G\ge1$, define
\begin{align}
 \cN_{m,n}(A,B,G;\cI)
 :={}&\#\bigl\{(j,r,s):j\in\cI,\ A<r\le2A,\ B<s\le2B,
 \notag\\
 &\hspace{7em}G<(r,s)\le2G,
 \ r\mid mj+h,\ s\mid nj+h\bigr\}.
 \label{eq:tpc29-incidence-count}
\end{align}
This is a triple count, not merely the number of occupied orbit
points.  At the left endpoint of a dyadic decomposition, the $G=1$
cell is understood as $1\le(r,s)\le2$; subsequent cells use the
displayed half-open convention.

\begin{theorem}[Fixed-row sparse incidence]
\label{thm:tpc29-incidence}
Uniformly for distinct physical rows $m,n$ and all $A,B,G\ge1$,
\begin{equation}
 \boxed{
 \cN_{m,n}(A,B,G;\cI)
 \ll_{\eps,h,\mathscr D}X^\eps
 \min\left\{J,\frac{AB}{G^2}+\frac JG\right\}.}
 \label{eq:tpc29-incidence-bound}
\end{equation}
Empty dyadic cells are understood to contribute zero.
\end{theorem}

\begin{proof}
First fix an exact common content $d=(r,s)$.  By
\cref{lem:tpc29-reflected-crt}, it must divide the nonzero integer
$h(m-n)$.  The number of possible exact $d$ in the dyadic interval is
therefore at most
\begin{equation}
 \tau(|h(m-n)|)\ll_{\eps,h,\mathscr D}X^\eps.
 \label{eq:tpc29-content-count}
\end{equation}
Write
\begin{equation}
 r=dr_1,\qquad s=ds_1,\qquad(r_1,s_1)=1.
 \label{eq:tpc29-primitive-parts}
\end{equation}
For a nonempty cell, $d\le2A$ and $d\le2B$.  Hence the number of
candidate pairs at this fixed $d$ is
\begin{equation}
 \ll\left(\frac Ad+1\right)
       \left(\frac Bd+1\right)
 \ll\frac{AB}{d^2}.
 \label{eq:tpc29-pair-count-fixed-d}
\end{equation}
Deleting the coprimality condition only enlarges the count.

Each compatible pair selects one residue class modulo
$q=rs/d$.  Since $r>A$ and $s>B$, its occupancy in $\cI$ is
\begin{equation}
 \ll_{\mathscr D}1+\frac Jq
 \ll_{\mathscr D}1+\frac{Jd}{AB}.
 \label{eq:tpc29-fiber-occupancy}
\end{equation}
Multiplying \eqref{eq:tpc29-pair-count-fixed-d} and
\eqref{eq:tpc29-fiber-occupancy}, using $d\asymp G$ (also in the
closed $G=1$ base cell), and then applying
\eqref{eq:tpc29-content-count} gives
\begin{equation}
 \cN_{m,n}(A,B,G;\cI)
 \ll X^\eps\left(\frac{AB}{G^2}+\frac JG\right).
 \label{eq:tpc29-incidence-crt-side}
\end{equation}

There is a second bound.  For each fixed $j$, the numbers of possible
$r$ and $s$ are at most $\tau(N_m(j))$ and $\tau(N_n(j))$.  The
pointwise divisor bound and $\#\cI\ll J$ give
\begin{equation}
 \cN_{m,n}(A,B,G;\cI)\ll_{\eps,\mathscr D}JX^\eps.
 \label{eq:tpc29-incidence-divisor-side}
\end{equation}
Taking the smaller of \eqref{eq:tpc29-incidence-crt-side} and
\eqref{eq:tpc29-incidence-divisor-side} proves the theorem.
\end{proof}

\begin{remark}[The indispensable $+1$]
\label{rem:tpc29-plus-one}
When $q\ge J$, \eqref{eq:tpc29-fiber-occupancy} is $O(1)$, not
$O(J/q)$.  Dropping this $+1$ would incorrectly turn every sparse
family into a saving.  The term $AB/G^2$ in
\eqref{eq:tpc29-incidence-bound} is precisely the total cost of those
one-point occupancies.
\end{remark}
